%% Springer Nature LaTeX template for The Journal of Supercomputing
%% https://link.springer.com/journal/11227/submission-guidelines
%% Converted from IEEEtran format

\documentclass[]{sn-jnl}

%% Packages
\usepackage[utf8]{inputenc}
\usepackage{amsmath,amssymb,amsfonts}
\usepackage{algorithmic}
\usepackage{graphicx}
\usepackage{textcomp}
\usepackage{xcolor}
\usepackage{listings}
\usepackage{booktabs}
\usepackage{array}
\usepackage{tabularx}
\usepackage{url}
\usepackage{enumitem}
\usepackage{colortbl}
\usepackage{multirow}
% Note: hyperref is already loaded by sn-jnl; do not load it again.

\setlist[itemize]{leftmargin=*, align=left}

\definecolor{codebackground}{rgb}{0.95,0.95,0.95}
\definecolor{tableheader}{rgb}{0.9,0.9,0.9}

\lstset{
    backgroundcolor=\color{codebackground},
    basicstyle=\footnotesize\ttfamily,
    breaklines=true,
    frame=single,
    numbers=left,
    numberstyle=\tiny,
    numbersep=6pt,
    xleftmargin=1.8em,
    framexleftmargin=1.8em,
    framesep=4pt,
    showstringspaces=false
}

%%------------------------------------------------------------
%%  TITLE AND AUTHORS
%%------------------------------------------------------------

\begin{document}

\title{PRM: Process Ancestry-Based Security for Defending HPC Environments from Privilege Escalation}

\author*[1]{\fnm{Mojtaba} \sur{Akbari}}
\email{mojtaba.akbari@gwdg.de}

\author*[1,2]{\fnm{Julian} \sur{Kunkel}}
\email{julian.kunkel@gwdg.de}

\affil*[1]{%
  \orgname{Gesellschaft f\"{u}r wissenschaftliche Datenverarbeitung mbH G\"{o}ttingen (GWDG)},
  \city{G\"{o}ttingen},
  \country{Germany}
}
\affil*[2]{%
  \orgname{Georg-August-Universit\"{a}t G\"{o}ttingen },
  \city{G\"{o}ttingen},
  \country{Germany}
}

%%------------------------------------------------------------
%%  ABSTRACT  (trimmed to ~143 words per Springer 100--150 word limit;
%%             one sentence removed, no text rephrased)
%%------------------------------------------------------------
\abstract{%
Modern High-Performance Computing (HPC) environments face a critical security
challenge: balancing robust protection with computational efficiency. Traditional
security approaches often introduce significant performance overhead or lack the
context awareness needed for complex HPC environments. This paper presents PRM
(Process Rule Management), a process ancestry-based security solution that
provides fine-grained control over system calls. By analyzing the lineage of
processes to make security decisions based on execution paths, PRM enables
administrators to distinguish between identical operations initiated through
different execution chains. PRM's intelligent caching achieves O(1) amortised performance
for repetitive operations, reducing redundant security checks in million-iteration
scientific computations. PRM can restrict even legitimate root administrators based
on execution context, enabling Trustworthy Research Environments (TRE) where
sensitive data analysis occurs with reduced reliance on traditional trust-based
security models. The eBPF-based architecture ensures system stability while
providing enhanced control over privileged operations in containerized HPC
environments.
}

\keywords{security, high-performance computing, process ancestry, system calls, containers, eBPF}

\maketitle

%%------------------------------------------------------------
%%  SECTION 1 — INTRODUCTION
%%------------------------------------------------------------
\section{Introduction}

Security in High-Performance Computing (HPC) environments presents unique
challenges that traditional security solutions struggle to address effectively.
The computational demands of scientific workloads require minimal security
overhead~\cite{zhang2021analyzing,xavier2013performance}, while the increasing
adoption of containerization introduces new attack vectors that must be
mitigated~\cite{reshetova2016security,rudyy2018containers}. Traditional security
mechanisms like SELinux, AppArmor, and seccomp lack the context awareness needed
to make nuanced security decisions in complex HPC
environments~\cite{provos2003improving,li2018speaker}, where identical system
calls may be legitimate or malicious depending on their execution context.

The ``privileged operation dilemma'' is particularly acute in HPC settings: how
to allow legitimate privileged operations required by scientific workflows while
preventing malicious exploitation~\cite{kurtzer2017singularity,priedhorsky2017charliecloud}.
Container technologies like Docker and Singularity often require root privileges,
creating significant security concerns in multi-tenant HPC clusters where
isolation between users and workloads is
paramount~\cite{canon2016scream,torrez2019hpc}.

\textbf{Research Question}: How can we develop a security framework that enables
safe execution of user-provided code in multi-tenant HPC environments while
maintaining the performance characteristics required for scientific computing?
This fundamental challenge requires distinguishing between legitimate scientific
computations and malicious activities that may appear identical at the system
call level~\cite{pappas2013kbouncer,ghavamnia2020confine}. Prior work has
demonstrated that syscall-level monitoring is fundamentally insufficient for
security decisions~\cite{wang2021semantic}, as HPC scientific workloads often
exhibit execution patterns that resemble malicious
behavior~\cite{vasilakis2019hpc}.

\textbf{Our Solution}: PRM addresses this challenge through process ancestry
analysis --- by examining not just \textit{``what''} is being executed, but
\textit{``how''} it came to be executed. Unlike traditional approaches that rely
on static syscall allowlists~\cite{ghavamnia2020confine}, PRM provides
context-aware privilege separation~\cite{ghavamnia2020confine} through semantic
execution context analysis. A legitimate scientific application launched through
Slurm $\rightarrow$ container $\rightarrow$ user process follows a predictable
ancestry pattern, while an exploit attempting privilege escalation exhibits
anomalous process lineage that PRM can detect and block. Our eBPF-based
implementation achieves low-overhead enforcement~\cite{nelson2020ebpf} while
providing unprecedented control over privileged operations, making possible
Trustworthy Research Environments where highly sensitive data can be analyzed
safely without relying on human trustworthiness.

\subsection{Contributions}

This paper makes the following contributions:

\begin{enumerate}
\item We introduce PRM, featuring dynamic rule redirection and runtime policy
      adaptation --- capabilities not available in existing static security solutions
\item We present an innovative multi-layer architecture combining intelligent
      caching, fast-decision filtering, and process ancestry analysis specifically
      designed for HPC environments
\item We demonstrate how PRM solves the trustable HPC environment challenge by
      enabling safe user code execution while preventing exploits through process
      ancestry validation
\item We provide implementation evidence showing how our caching mechanism
      eliminates redundant security checks for repetitive operations (e.g.,
      million-iteration syscalls like file writes, file opens, and network operations)
\end{enumerate}

%%------------------------------------------------------------
%%  SECTION 2 — RELATED WORK
%%------------------------------------------------------------
\section{Related Work}
\label{sec:related}

\subsection{System Call Filtering Approaches}

System call filtering has been explored extensively, progressing from static
policies to programmable enforcement.
Systrace~\cite{provos2003improving} was among the earliest systems to enforce
per-process system call policies, but its user-space policy engine introduced
context-switch overhead that is prohibitive for high-throughput HPC workloads.
Seccomp-BPF~\cite{edge2015seccomp} moved filtering into the kernel via cBPF
programs attached at process level; however, its policies are fixed at load time
and cannot incorporate runtime context such as process ancestry or UID
transitions.
KRSI~\cite{singh2019krsi} extends the eBPF-based model to LSM hooks, providing
programmable security instrumentation, yet it offers no built-in ancestry
tracking, caching, or dynamic rule redirection --- administrators must implement
these from scratch.
Landlock~\cite{salvan2016landlock} enables unprivileged, stackable access
control through eBPF, but is scoped to filesystem and network restrictions and
cannot express policies that depend on how a process was spawned.

PRM builds on the programmability introduced by KRSI and Landlock while adding
three capabilities absent from all of the above: multi-generational ancestry
analysis, intelligent per-syscall caching, and dynamic rule redirection at
runtime.

\subsection{Container Security Solutions}

Container isolation has been addressed at different levels of the software stack.
gVisor~\cite{google2018gvisor} interposes a user-space kernel that re-implements
Linux system calls, providing strong isolation at the cost of significant
performance overhead and incomplete syscall coverage --- both of which conflict
with HPC requirements for native performance and full POSIX
semantics~\cite{torrez2019hpc}.
Kata Containers~\cite{kata2018architecture} achieve isolation through lightweight
VMs, adding hardware-level boundaries but introducing memory and boot-time
overhead that limits suitability for short-lived, high-throughput HPC
jobs~\cite{xavier2013performance}.
Falco~\cite{falco2016runtime} monitors container runtime behaviour via kernel
event streams and raises alerts, but operates in a detection-only mode; it cannot
enforce policy or block operations inline.

PRM complements these approaches by operating inside the kernel's LSM path,
enabling inline enforcement with process-ancestry context --- a combination that
allows it to block container escape attempts at the syscall level rather than
detecting them after the fact.

\subsection{Process Ancestry and Execution-Context Research}

Prior work has explored process relationships and execution paths for security
decisions, though with narrower scope than PRM.
YAMA~\cite{cook2012yama} restricts \texttt{ptrace} attachment based on
parent--child relationships, effectively a single-generation ancestry check
limited to one LSM hook.
SPEAKER~\cite{li2018speaker} analyses execution paths to partition application
kernels into security phases, but targets intra-process control flow rather than
inter-process lineage and does not provide a general-purpose policy framework.
Confine~\cite{ghavamnia2020confine} automatically generates per-container seccomp
policies by static analysis of container images; its policies are fixed at
container start and cannot adapt to runtime ancestry or UID transitions.
Wang et al.~\cite{wang2021semantic} propose semantic-aware monitoring for
containerised applications, demonstrating that syscall-level observation alone is
insufficient and that higher-level context improves detection accuracy.

PRM extends these insights into a unified enforcement framework: it combines
multi-generational ancestry tracking (configurable depth, currently 3 levels),
UID transition vector analysis across the ancestry chain, dynamic rule
redirection at runtime, and specialised PROG handlers per LSM hook type ---
capabilities that, taken together, are not available in any single existing
system.

%%------------------------------------------------------------
%%  SECTION 3 — PRM FRAMEWORK DESIGN
%%------------------------------------------------------------
\section{PRM Framework Design}
\label{sec:design}

\subsection{Architecture Overview}

Figure~\ref{fig:architecture} illustrates PRM's multi-layer architecture designed
for HPC environments where performance and security must coexist.

\begin{figure}[htbp]
\begin{lstlisting}[basicstyle=\tiny\ttfamily, frame=single, breaklines=true, language=c]
PRM_SECURITY_CHECK(syscall_context) {
  // Layer 1: Intelligent Caching
  key = GENERATE_KEY(pid, tid, hook_type)
  if (cached_decision = CACHE_LOOKUP(key)) 
    return cached_decision  // O(1) performance
  
  // Layer 2: Fast-Decision Filters  
  filter_result = APPLY_FILTERS(syscall_context)
  if (filter_result == IMMEDIATE_ALLOW) {
    CACHE_UPDATE(key, ALLOW)
    return ALLOW
  }
  if (filter_result == IMMEDIATE_DENY) {
    CACHE_UPDATE(key, DENY) 
    return DENY
  }
  
  // Layer 3: Process Ancestry Analysis
  current = GET_PROCESS_NAME(syscall_context.task)
  parent = GET_PARENT_NAME(syscall_context.task_parent)
  grandparent = GET_GRANDPARENT_NAME(syscall_context.task_grandparent)
  
  // Layer 4: Dynamic Rule Processing
  for (i = 0; i < MAX_RULES; i++) {
    rule = PRM_TABLE[i]
    if (MATCH(current, rule.process) && 
        MATCH(parent, rule.parent) && 
        MATCH(grandparent, rule.grandparent)) {
      
      if (rule.action == REDIRECT) {
        i = rule.redirect_index  // Dynamic rule chaining
        continue
      }
      if (rule.action == RETURN) {
        result = PROG_HANDLER[rule.prog_num](syscall_context)
        CACHE_UPDATE(key, result)
        return result
      }
      CACHE_UPDATE(key, rule.action)
      return rule.action
    }
  }
  
  // Layer 5: Default Policy
  CACHE_UPDATE(key, ALLOW)
  return ALLOW  // Fail-safe default
}
\end{lstlisting}
\caption{PRM Architecture: Five-layer security processing with intelligent
caching, fast filtering, ancestry analysis, dynamic rules, and fail-safe
defaults.}
\label{fig:architecture}
\end{figure}

PRM implements a five-layer architecture:

\begin{itemize}
\item \textbf{Layer 1 --- Intelligent Caching}: O(1) lookup for previously
      evaluated syscalls
\item \textbf{Layer 2 --- Fast-Decision Filtering}: Immediate allow/deny for
      common patterns
\item \textbf{Layer 3 --- Process Ancestry Analysis}: Multi-generational lineage
      extraction
\item \textbf{Layer 4 --- Dynamic Rule Processing}: Runtime rule redirection and
      PROG handlers
\item \textbf{Layer 5 --- Default Policy}: Fail-safe allow with comprehensive
      logging
\end{itemize}

\subsection{Key Innovations}

\subsubsection{Ancestry-Based Security: Beyond Traditional Privilege Models}

PRM introduces an approach to system security that extends traditional
privilege-based models. Whereas conventional security frameworks operate on the
principle of ``who you are'' (user identity) or ``what you want to do''
(capability-based), PRM adds a third dimension: ``how you got here'' (process
ancestry). Prior work has explored aspects of this idea --- YAMA~\cite{cook2012yama}
restricts \texttt{ptrace} based on parent--child relationships, and
SPEAKER~\cite{li2018speaker} analyses execution paths for intra-process security
--- but neither provides a general-purpose, multi-generational enforcement
framework that spans arbitrary LSM hooks. PRM combines and extends these concepts
into a unified policy engine, enabling the following capabilities:

\textbf{Restricting Root Privileges Through Ancestry Analysis}: PRM can restrict
even root operations based on execution context. A root process spawned through an
unexpected ancestry chain (e.g., compromised container $\rightarrow$ privilege
escalation $\rightarrow$ root shell) can be denied critical operations, while the
same root process launched through legitimate paths (e.g., systemd $\rightarrow$
slurmstepd $\rightarrow$ container $\rightarrow$ root) operates normally. This
makes privilege escalation attacks less effective even when the attacker obtains
root credentials.

\textbf{Context-Aware Privilege Enforcement}: Traditional security models provide
limited protection once attackers gain legitimate credentials. PRM's
ancestry-based approach means that even with valid root credentials, processes are
constrained by their execution lineage. This creates an ``execution context
firewall'' that limits lateral movement and privilege abuse regardless of
credential compromise.

\textbf{Trustable Environment Creation}: By analysing process ancestry, PRM
enables the creation of computing environments where even privileged operations
are contextually validated. This allows HPC centres to execute untrusted user
code while maintaining system integrity.

\textbf{Trustworthy Research Environment (TRE) Implementation}: PRM's
ancestry-based approach enables the creation of Trustworthy Research Environments
for highly sensitive data analysis. Unlike traditional security models that rely
on user trust, PRM can restrict even legitimate administrators and root users based
on their execution context. In a TRE deployment, administrators with valid root
credentials are constrained to specific operational patterns --- they can perform
system maintenance through approved paths (e.g., SSH $\rightarrow$ systemd
$\rightarrow$ maintenance scripts) but are blocked from accessing sensitive research
data through unauthorised lineages (e.g., SSH $\rightarrow$ shell $\rightarrow$
data access). This provides a level of data protection where even system
administrators cannot bypass security controls, enabling safe remote analysis of
health records, financial data, and classified research without data exfiltration
risks. The ancestry-based enforcement ensures that sensitive data remains within the
controlled environment regardless of privilege level or credential compromise.

\subsubsection{Dynamic Rule Redirection}

Unlike static security solutions that rely on fixed policies, PRM introduces
dynamic rule redirection capabilities that adapt at runtime. The framework
supports two primary redirection mechanisms: REDIRECT actions that can jump to
any of the 300 available rule indices, and RETURN actions that invoke specialized
PROG handlers for complex decision logic. This creates a hierarchical rule
organization with protected zones for different security contexts, enabling
runtime policy adaptation based on execution environment. Such dynamic behavior
is impossible with traditional static approaches like SELinux or AppArmor.

\subsubsection{Process Ancestry Analysis}

PRM's ancestry analysis provides multi-generational process lineage tracking:

\begin{table}[htbp]
\centering
\caption{Process Ancestry Analysis Components}
\label{tab:ancestry}
\begin{tabular}{@{}p{2.5cm}p{5cm}@{}}
\toprule
\textbf{Component} & \textbf{Description} \\
\midrule
Current Process   & Active process making system call \\
Parent Process    & Direct parent in process tree \\
Grandparent Process & Second-level ancestor \\
UID Transitions   & User ID changes across ancestry \\
Execution Path    & Complete lineage validation \\
HookType          & LSM hook identifier for syscall classification \\
HookStructure     & Data structure containing syscall context \\
\bottomrule
\end{tabular}
\end{table}

\subsubsection{Worst-Case Performance for Each Pattern}

PRM's performance characteristics vary by execution pattern, with worst-case
scenarios defined for each processing layer:

\begin{table}[htbp]
\centering
\caption{PRM Worst-Case Performance Analysis}
\label{tab:performance}
\begin{tabular}{@{}p{2cm}p{3cm}p{2.5cm}@{}}
\toprule
\textbf{Layer} & \textbf{Function} & \textbf{Worst-Case} \\
\midrule
Cache Lookup     & Decision retrieval  & O(1) hash collision \\
Fast Filters     & Immediate decisions & O(n) filter checks \\
Ancestry Analysis & Lineage extraction & O(depth) traversal \\
Rule Processing  & Policy evaluation   & O(MAX\_RULES) scan \\
\bottomrule
\end{tabular}
\end{table}

\subsubsection{Intelligent Caching for HPC Workloads}

HPC applications often perform repetitive operations (e.g., million-iteration file
writes). First syscall extracts three critical data parameters (p1, p2, p3) from
the hook structure, where each parameter represents context-specific information
converted to integer values for efficient processing. These parameters create a
unique context-aware syscall signature that enables precise cache differentiation.
Due to kernel space constraints that prohibit complex hashing algorithms, PRM
employs the Jenkins hash function for its speed and maintainability:
hash = Jenkins(p1, p2, p3, PID, TID, GID, hook\_type). Subsequent identical
syscalls achieve O(1) lookup with total overhead:
$1 \times T_{\mathrm{security}} + 999{,}999 \times T_{\mathrm{cache}}$
where $T_{\mathrm{cache}} \ll T_{\mathrm{security}}$.

Our caching system:

\begin{itemize}
\item \textbf{UniqueKey Generation}: Combines PID, Thread ID, and Hook Type
\item \textbf{LRU Eviction}: Maintains 2024 most recent decisions
\item \textbf{Hook Validation}: Prevents security bypasses through cache poisoning
\item \textbf{Performance Impact}: First syscall pays full cost, subsequent
      identical calls are O(1)
\end{itemize}

\textbf{Example}: A loop writing to the same file 1,000,000 times:
\begin{itemize}
\item Traditional systems: 1,000,000 full security checks
\item PRM Framework: 1 full check + 999,999 cache hits
\end{itemize}

\subsubsection{Fast-Decision Filtering}

Filters provide immediate decisions before expensive ancestry analysis:

\begin{itemize}
\item \textbf{Self-Process Detection}: Immediate allow for framework's own operations
\item \textbf{Known-Good Patterns}: Fast-path for common legitimate operations
\item \textbf{Early Rejection}: Block obviously malicious calls without full processing
\item \textbf{Kernel Syscall Bypass}: Configurable bypass for kernel and kernel
      thread operations
\item \textbf{Context Enrichment}: Prepare data for downstream components
\end{itemize}

PRM encourages bypassing kernel syscalls to achieve significant performance
improvements in HPC environments. While kernel syscalls can be analyzed for
security threats, they are typically trusted operations originating from the
kernel itself. Bypassing these calls eliminates unnecessary overhead during
intensive scientific computations where kernel threads frequently perform memory
management, scheduling, and I/O operations. This optimization is particularly
beneficial for HPC workloads that generate high-frequency kernel activity.

\subsection{HPC User Safety Model}

The PRM Framework addresses a critical HPC challenge: enabling users to safely
execute custom code and instructions while preventing system compromise.

\subsubsection{Safe Code Execution Environment}

HPC centers must allow users to:
\begin{itemize}
\item Install and run custom scientific software
\item Execute containerized applications with elevated privileges
\item Access shared filesystems and compute resources
\item Perform legitimate system operations required by scientific workflows
\end{itemize}

Traditional security solutions are inadequate because they:
\begin{itemize}
\item Lack context to distinguish legitimate vs.\ malicious privileged operations
\item Introduce unacceptable performance overhead for compute-intensive workloads
\item Cannot adapt to the dynamic nature of scientific computing environments
\item Fail to provide fine-grained control over container security
\end{itemize}

\subsubsection{PRM's Approach to User Safety}

PRM enables safe user code execution through:

\begin{enumerate}
\item \textbf{Process Ancestry Validation}: Ensures code executes through
      expected paths (e.g., Slurm --- container --- user application)
\item \textbf{UID Transition Monitoring}: Detects privilege escalation using
      vector-based pattern matching (19 patterns including [0,0,0,0] for root
      daemons, [0,2,2,2] for escalation, [2,2,2,2] for users, [0,2,2,3] for
      critical mismatches)
\item \textbf{Container Escape Prevention}: Analyzes execution context to prevent
      container breakouts
\item \textbf{Dynamic Policy Adaptation}: Adjusts security policies based on
      execution environment
\end{enumerate}

\subsection{HPC Profile Rules}

PRM implements specialized rule configurations for different HPC deployment
scenarios through a 300-entry rule table with dynamic redirection capabilities.
Table~\ref{tab:hpc_profiles} shows key rule patterns extracted from our reference
configuration designed for HPC environments, while Table~\ref{tab:prm_config}
details the comprehensive configuration parameters that govern PRM's behavior in
HPC environments.

\begin{table}[htbp]
\centering
\caption{HPC Profile Rule Configuration}
\label{tab:hpc_profiles}
\begin{tabular}{@{}p{1.8cm}p{2.2cm}p{1.5cm}p{1.2cm}p{0.8cm}@{}}
\toprule
\textbf{Process} & \textbf{Parent} & \textbf{Action} & \textbf{Hook} & \textbf{Zone} \\
\midrule
\multicolumn{5}{@{}l@{}}{\textit{System Services (Rules 1--15)}} \\
systemd    & *              & ACCEPT   & ALL      & 0 \\
runc       & containerd-shim & ACCEPT  & ALL      & 0 \\
slurmstepd & systemd        & ACCEPT   & ALL      & 0 \\
\midrule
\multicolumn{5}{@{}l@{}}{\textit{Blocked Interpreters (Rules 25--42)}} \\
ruby & * & REJECT & ALL & 0 \\
node & * & REJECT & ALL & 0 \\
gdb  & * & REJECT & ALL & 0 \\
\midrule
\multicolumn{5}{@{}l@{}}{\textit{Hook Redirections (Rules 43--53)}} \\
* & * & REDIRECT & KILL\_SIG  & 290 \\
* & * & REDIRECT & FILE\_OPEN & 62  \\
* & * & REDIRECT & CAPABLE    & 62  \\
\midrule
\multicolumn{5}{@{}l@{}}{\textit{Protected Pattern-1 (Rules 62--74)}} \\
*      & slurmstepd & RETURN & ALL & 1 \\
*      & runc       & RETURN & ALL & 1 \\
python & *          & RETURN & ALL & 1 \\
\ldots & \ldots     & \ldots & \ldots & \ldots \\
\bottomrule
\end{tabular}
\end{table}

\begin{table}[htbp]
\centering
\caption{PRM Framework Configuration Parameters}
\label{tab:prm_config}
\begin{tabular}{@{}p{2.5cm}p{2.8cm}p{2.2cm}@{}}
\toprule
\textbf{Category} & \textbf{Parameter} & \textbf{Values} \\
\midrule
\multicolumn{3}{@{}l@{}}{\textit{User/Group Security}} \\
Safe UIDs  & System Operations  & 0, 1000, 1001 \\
Unsafe GIDs & Restricted Groups & 0, 1000, 2000 \\
\midrule
\multicolumn{3}{@{}l@{}}{\textit{Directory Access Control}} \\
Valid Directories & HPC Data Paths    & home, scratch, work \\
Blocked Paths     & Dangerous Access  & /proc/1/mem, /proc/kcore \\
Binary Directories & System Binaries  & /sbin, /bin, /usr \\
\midrule
\multicolumn{3}{@{}l@{}}{\textit{Network Security}} \\
Socket Protocols & Blocked Protocols   & AF\_PACKET, AF\_IB \\
IP Destination   & Restricted Networks & 169.254.0.0/16 \\
IP Source        & Allowed Networks    & 10.0.0.0/8 \\
\midrule
\multicolumn{3}{@{}l@{}}{\textit{Execution Control}} \\
BPRM Destinations & Blocked Exec Paths  & /dev/shm/, /tmp/. \\
Interpreters      & Restricted Binaries & perl, R, singularity \\
Module Deny       & Blocked Modules     & ib\_, nvidia, kvm, \ldots \\
\bottomrule
\end{tabular}
\end{table}

\subsection{PROG Security Handlers}

PRM implements specialized PROG handlers that provide hook-specific security logic
when invoked through RETURN actions. Our comprehensive testing framework validates
protection against 13 major attack vectors commonly targeting HPC environments.
Table~\ref{tab:prog_handlers} details the complete set of PROG security handlers
and their specific functions.

\begin{table}[htbp]
\centering
\caption{PRM PROG Security Handlers}
\label{tab:prog_handlers}
\begin{tabular}{@{}p{1cm}p{5cm}p{4cm}@{}}
\toprule
\textbf{ID} & \textbf{Handler Name} & \textbf{Security Function} \\
\midrule
1     & signalKillTracer              & Prevents process termination attacks \\
2     & denyWriteOutOfValidDirectories & Restricts file writes to safe paths \\
3     & denyMakeSocketToEndHost       & Blocks dangerous network connections \\
4     & bprmSecurityCheck             & Validates binary execution paths \\
5     & memoryProtectCheck            & Prevents code injection via mprotect \\
6     & denyIncomeSocket              & Controls incoming socket connections \\
7     & denyFileOpen                  & Restricts file access based on ancestry \\
8     & denyLoadModule                & Blocks dangerous kernel module loading \\
9     & credPrepareCheck              & Validates credential transitions \\
10    & capableCheck                  & Controls capability usage \\
11    & taskFixSetUIDCheck            & Prevents setuid exploitation \\
12--31 & EMPTY\_PROG                 & Reserved for future extensions \\
\bottomrule
\end{tabular}
\end{table}

PRM addresses critical security threats through comprehensive attack vector
coverage and specialized PROG Security Handlers. These handlers provide
hook-specific security logic tailored to different LSM interception points,
enabling fine-grained control over file operations, network connections,
credential management, and capability checks. Table~\ref{tab:attack_coverage}
demonstrates our comprehensive coverage of attack vectors with recent real-world
examples, showing the relevance and effectiveness of PRM's security protection.

\begin{table*}[t]
\centering
\caption{Attack Vector Prevalence and PRM Protection Coverage (2024--2025)}
\label{tab:attack_coverage}

\resizebox{\textwidth}{!}{%
\begin{tabular}{lllll}
\toprule
\textbf{Attack Vector} & \textbf{Public Category} & \textbf{Prevalence} & \textbf{PROG} & \textbf{References} \\
\midrule
Network Reconnaissance  & Pre-Attack Scanning      & \textbf{41\%} & PROG\_3,6  & Verizon DBIR 2024~\cite{verizon2024dbir} \\
Memory Code Injection   & Memory Corruption RCE    & \textbf{32\%} & PROG\_5    & ENISA TL2024~\cite{enisa2024threat} \\
Interpreter Abuse       & Scripting Engine Abuse   & \textbf{29\%} & PROG\_4    & CrowdStrike GTR 2024~\cite{crowdstrike2024gtr} \\
Credential Manipulation & Credential Theft         & \textbf{29\%} & PROG\_9    & Verizon DBIR 2024~\cite{verizon2024dbir} \\
Sensitive File Access   & Unauthorized File Access & \textbf{22\%} & PROG\_7    & ENISA TL2024~\cite{enisa2024threat} \\
Process Memory Access   & Credential Dumping       & \textbf{17\%} & PROG\_7    & MITRE ATT\&CK 2024~\cite{mitre2024attack} \\
Privilege Escalation    & Priv-Esc Vulnerabilities & \textbf{12\%} & PROG\_11   & Verizon DBIR 2024~\cite{verizon2024dbir} \\
Socket Binding Attacks  & Port Hijacking           & \textbf{10\%} & PROG\_6    & ENISA TL2024~\cite{enisa2024threat} \\
Container Escape        & Container Intrusions     & \textbf{9\%}  & PROG\_4,7  & CrowdStrike GTR 2024~\cite{crowdstrike2024gtr} \\
Directory Traversal     & Path Traversal Abuse     & \textbf{8\%}  & PROG\_2    & OWASP Top 10 2024~\cite{owasp2024top10} \\
Capability Escalation   & Linux Capabilities       & \textbf{7\%}  & PROG\_10   & Wiz Cloud Report 2024~\cite{wiz2024cloud} \\
Process Termination     & Process Manipulation     & \textbf{5\%}  & PROG\_1    & CrowdStrike GTR 2024~\cite{crowdstrike2024gtr} \\
Kernel Module Loading   & Kernel-Level Implants    & \textbf{4\%}  & PROG\_8    & Mandiant M-Trends 2024~\cite{mandiant2024mtrends} \\
\bottomrule
\end{tabular}%
}
\end{table*}

The global vulnerability exploitation landscape (2023--2025) provides critical
context for PRM's security coverage. Recent industry reports show 768 CVEs were
exploited in 2024 (up from 639 in 2023), with 23.6\% being zero-day exploits and
vulnerability exploitation accounting for 14\% of all breaches --- representing a
180\% increase compared to
2023~\cite{vulncheck2024trends,verizon2024dbir}. These data reflect publicly
reported vulnerability-exploitation events globally across all sectors. Because
public breach reporting aggregates many exploit types under `vulnerability
exploitation' and generally does not reveal whether the target was an HPC/data-center,
the statistics should be considered a lower bound on exposure. Finer-grained
breakdowns (container escape, kernel rootkits, memory-code injection, etc.) are
not systematically reported; their prevalence in HPC environments remains unknown.

PRM's PROG handlers are designed to address the most critical threats facing
modern computing environments based on current attack vector prevalence data.

%%------------------------------------------------------------
%%  SECTION 4 — EVALUATION
%%------------------------------------------------------------
\section{Evaluation}
\label{sec:evaluation}

\subsection{Implementation Validation}

We have successfully implemented and validated the PRM Framework through
comprehensive testing and laboratory evaluation. The complete implementation
comprises 6,249 lines of code across 39 source files: 5,281 LOC in C/header
files (2,558 LOC in C implementation files and 2,723 LOC in header files) plus
968 LOC in supporting Python, shell scripts, and other utilities, demonstrating
substantial engineering effort and comprehensive functionality coverage. The
implementation demonstrates complete integration with the Linux Security Module
framework, supporting 12 specialized PROG handlers across different LSM hooks
including FILE\_OPEN, SOCKET\_CONNECT, CRED\_PREPARE, and CAPABLE operations.
The framework has been thoroughly tested and validated in controlled laboratory
environments.

\subsubsection{Functional Validation}

The rule engine operates with full 300-entry capacity, successfully processing
complex ancestry patterns and dynamic redirections. Our caching system demonstrates
effective LRU eviction behavior, maintaining consistent O(1) lookup performance
across varying workload patterns. The UID transition analysis component correctly
identifies all 19 pre-trained attack patterns, with successful detection rates
validated through controlled privilege escalation scenarios.

\subsubsection{Unique Security Capabilities}

Table~\ref{tab:capabilities} shows security capabilities where PRM provides
unique functionality not available in existing solutions:

\begin{table}[htbp]
\caption{Security capabilities comparison}
\label{tab:capabilities}
\centering
\begin{tabular}{@{}p{3.0cm}p{0.3cm}p{0.7cm}p{0.9cm}p{0.8cm}@{}}
\toprule
\textbf{Capability} & \textbf{PRM} & \textbf{SELinux} & \textbf{AppArmor} & \textbf{Seccomp} \\
\midrule
Process ancestry analysis       & Yes & No      & No      & No \\
Programmable Policy Extensions  & Yes & No      & No      & No \\
UID transition detection        & Yes & Limited & Limited & Limited \\
Container escape prevention     & Yes & Yes     & Limited & Limited \\
Multi-generational tracking     & Yes & No      & No      & No \\
Dynamic rule redirection        & Yes & Limited & Limited & Limited \\
\bottomrule
\end{tabular}
\end{table}

\subsection{Performance Measurements}
\label{sec:performance_measurements}

We conducted performance analysis on our implementation to validate the
theoretical advantages of our multi-layer architecture.

\subsubsection{Cache Effectiveness Measurements}

Testing with a synthetic HPC workload performing 1,000,000 file write operations
showed dramatic performance improvements. Traditional security frameworks require
full policy evaluation for each system call, resulting in linear overhead scaling.
PRM's caching mechanism reduces this to a single full evaluation plus 999,999
O(1) cache lookups, demonstrating the practical impact of our intelligent caching
design.

\subsubsection{Memory Footprint Analysis}

Our implementation maintains a consistent memory footprint of approximately 4--5
MB per compute node, measured across different workload scenarios. The LRU cache
with 2024 entries provides optimal balance between memory usage and hit rates,
while the 300-entry rule table and 19-pattern UID database contribute minimal
overhead. This footprint remains constant regardless of the number of concurrent
processes, making PRM suitable for large-scale HPC deployments.

\subsubsection{Expected Performance Characteristics}

Based on eBPF performance literature~\cite{hoiland2018express} and our
architecture:
\begin{itemize}
\item \textbf{Cache Hit Scenario}: Near-zero overhead (hash lookup only)
\item \textbf{Cache Miss Scenario}: Bounded processing time due to fixed rule
      table size
\item \textbf{Memory Footprint}: Approximately 4--5 MB per node (measured)
\item \textbf{Scalability}: Linear scaling (per-node processing)
\item \textbf{System Safety}: eBPF verifier guarantees prevent kernel crashes and
      ensure bounded execution time~\cite{nelson2017bpf}, making PRM safe for
      production HPC environments
\end{itemize}

\subsection{HPC-Specific Benefits}

The PRM Framework addresses critical HPC security challenges:

\subsubsection{Multi-Tenant Security}

HPC centers serve multiple users and projects simultaneously:
\begin{itemize}
\item \textbf{User Isolation}: Process ancestry ensures jobs don't escape user
      boundaries
\item \textbf{Resource Protection}: Prevents unauthorized access to other users'
      data
\item \textbf{Privilege Containment}: Limits scope of privileged operations
\end{itemize}

\subsubsection{Container Security for Scientific Computing}

Scientific workflows increasingly rely on containers:
\begin{itemize}
\item \textbf{Singularity Integration}: Secure execution of scientific containers
\item \textbf{Docker Support}: Safe privileged container operations
\item \textbf{Escape Prevention}: Process ancestry detects container breakout attempts
\end{itemize}

\subsubsection{Workload Manager Integration}

Compatibility with HPC schedulers:
\begin{itemize}
\item \textbf{Slurm Integration}: Rules for slurmstepd and slurmd processes
\item \textbf{PBS/Torque Support}: Ancestry-based job isolation
\item \textbf{Custom Schedulers}: Flexible rule system adapts to any scheduler
\end{itemize}

\subsection{Security Analysis}

We analyzed our framework's security effectiveness through code review and
theoretical analysis:

\subsubsection{Attack Vector Coverage}

Our implementation addresses:
\begin{enumerate}
\item \textbf{Container Escapes}: Process ancestry rules detect unusual execution
      paths from containers
\item \textbf{Privilege Escalation}: UID transition analysis identifies 19 known
      attack patterns including:
   \begin{itemize}
   \item Direct root compromise (severity: CRITICAL)
   \item UID mismatch chains (severity: CRITICAL)
   \item Setuid exploitation (severity: ESCALATED)
   \item Alternating privilege patterns (severity: ANOMALOUS)
   \end{itemize}
\item \textbf{Capability Abuse}: CAPABLE hook blocks dangerous capabilities:
   \begin{itemize}
   \item CAP\_SETUID: Prevents setuid() attacks
   \item CAP\_SETGID: Prevents setgid() attacks
   \item CAP\_SYS\_ADMIN: Restricted based on lineage analysis
   \item CAP\_DAC\_OVERRIDE: Limited to root processes only
   \end{itemize}
\item \textbf{File System Attacks}: Directory-based rules prevent unauthorized
      access patterns
\item \textbf{Network Attacks}: Socket rules control connection establishment
      based on ancestry
\item \textbf{Memory Attacks}: MPROTECT hook prevents code injection via
      executable memory
\end{enumerate}

\subsubsection{Dynamic Security Adaptation}

PRM's dynamic rule system enables:
\begin{itemize}
\item \textbf{Context-Aware Policies}: Different rules for different execution
      contexts
\item \textbf{Runtime Rule Chaining}: REDIRECT actions create complex policy flows
\item \textbf{Specialized Handlers}: PROG functions provide hook-specific security
      logic
\item \textbf{Protected Zones}: Hierarchical rule organization (e.g., container
      zones, HPC zones)
\end{itemize}

\subsubsection{Security Model Validation}

Our security model provides:
\begin{itemize}
\item \textbf{Defense in Depth}: Multiple layers (cache, filters, verifier, PROGs)
\item \textbf{Fail-Safe Defaults}: Unknown patterns default to ACCEPT with logging
\item \textbf{Audit Trail}: Comprehensive logging of all security decisions
\item \textbf{Policy Validation}: Rule syntax checking and circular reference
      detection
\item \textbf{Cache Security}: Hook-specific validation prevents cache poisoning
\end{itemize}

\subsubsection{UID Ancestry Analysis Algorithm}

PRM implements a sophisticated window-sliding algorithm for detecting malicious
UID transition vectors in process ancestry chains. Due to eBPF kernel space
limitations that prevent complex machine learning algorithms like k-NN or
clustering, we developed an efficient pattern matching approach using sliding
window analysis.

\begin{lstlisting}[
    caption={UID Ancestry Analysis Algorithm},
    captionpos=t,
    framesep=8pt,
    belowcaptionskip=5pt,
    language=C++
]
// UID Ancestry Analysis Algorithm
// Input: Current process task, unique key
// Output: Security severity level

int build_ancestry_chain(task_struct* task, uint64_t key) {
  int ancestry_buffer[MAX_ANCESTORS];
  uid_t current_uid = get_uid(task);
  ancestry[0] = encode_uid(current_uid); // 0=ROOT, 2=USER, 3=MISMATCH
  
  task_struct* parent_process = task->parent;
  for (int i = 1; i < MAX_ANCESTORS; i++) {
    if (parent_process == NULL) {
      ancestry[i] = EMPTY_CELL;
    } else {
      uid_t parent_uid = get_uid(parent_process);
      ancestry[i] = detect_transition(current_uid, parent_uid);
      parent_process = parent_process->parent;
    }
  }
  return sliding_window_match(ancestry);
}

int sliding_window_match(int ancestry[]) {
  int max_severity = 0;
  
  for (int p = 0; p < ATTACK_VECTOR_COUNT; p++) {
    pattern_t pattern = ATTACK_VECTOR_DATABASE[p];
    for (int window_pos = 0; window_pos <= 3; window_pos++) {
      int offset = window_pos * 4;
      if (match_pattern(&ancestry[offset], pattern)) {
        if (pattern.severity > max_severity) {
          max_severity = pattern.severity;
        }
      }
    }
  }
  return max_severity;
}
\end{lstlisting}

The algorithm operates in two phases: ancestry chain construction with symbolic
UID encoding, followed by sliding window pattern matching against 19 pre-trained
attack vectors. This approach circumvents eBPF limitations while providing
efficient malicious pattern detection with O(n*m) complexity where n=ancestry
depth and m=pattern database size.

\subsection{Deployment Readiness and Production Integration}

\subsubsection{HPC Center Integration}

Integrating PRM into production HPC environments will require:
\begin{itemize}
\item \textbf{Scheduler Integration}: Rules must accommodate Slurm, PBS, and
      custom schedulers
\item \textbf{User Training}: Administrators need training on ancestry-based
      security concepts
\item \textbf{Policy Development}: Site-specific rules based on local workflows
      and requirements
\item \textbf{Monitoring Integration}: Connection to existing security monitoring
      infrastructure
\end{itemize}

\subsubsection{Performance Considerations}

For HPC deployments:
\begin{itemize}
\item \textbf{Cache Sizing}: Larger caches for systems with many concurrent jobs
\item \textbf{Rule Optimization}: Frequently-matched rules should appear earlier
      in table
\item \textbf{Filter Tuning}: Site-specific filters for common application patterns
\item \textbf{Memory Management}: Efficient buffer allocation for high-throughput
      systems
\end{itemize}

\subsection{Security Validation Through Attack Simulation}
\label{sec:security_validation}

We evaluated PRM's security capabilities through controlled laboratory testing of
attack scenarios. These scenarios illustrate PRM's approach to correlating process
ancestry with security decisions.

\subsubsection{Setuid Binary Exploitation}

\textbf{Attack}: Setuid Binary Exploitation
\begin{itemize}
\item Attacker exploits a vulnerable setuid binary inside a container
\item The binary spawns a shell with elevated privileges
\item Expected lineage: \texttt{container $\rightarrow$ application $\rightarrow$ worker}
\item Observed lineage: \texttt{container $\rightarrow$ setuid\_binary $\rightarrow$ root\_shell $\rightarrow$ sensitive\_files}
\item PRM detects an unexpected UID transition [1000 $\rightarrow$ 0] within the
      process ancestry
\item Existing MAC systems (e.g., SELinux, AppArmor) permit setuid execution under
      policy, but do not model or detect anomalous privilege transitions based on
      runtime ancestry
\end{itemize}

\subsubsection{Unauthorized Shared Directory Access}

\textbf{Attack}: Unauthorized Shared Directory Access
\begin{itemize}
\item Multiple users share a \texttt{/project} directory for collaboration
\item User A's job attempts to modify User B's private files
\item Expected lineage: \texttt{slurm $\rightarrow$ userA\_job $\rightarrow$ userA\_files}
\item Observed lineage: \texttt{slurm $\rightarrow$ userA\_job $\rightarrow$ userB\_private\_files}
\item PRM can block operations by correlating process ancestry with target file
      ownership
\item Existing MAC systems (e.g., SELinux, AppArmor) enforce access based on
      current credentials and static labels, but lack ancestry-aware policies that
      bind job provenance to permissible file ownership
\end{itemize}

\subsubsection{Malicious Process Injection via Shared Directory}

\textbf{Attack}: Malicious Process Injection via Shared Directory
\begin{itemize}
\item Attacker places a malicious script in a shared \texttt{/scratch} directory
\item The script is executed by a legitimate HPC job
\item The injected process attempts to access other users' data via the shared
      filesystem
\item Expected lineage: \texttt{slurm $\rightarrow$ user\_job $\rightarrow$ trusted\_script}
\item Observed lineage: \texttt{slurm $\rightarrow$ user\_job $\rightarrow$ malicious\_script $\rightarrow$ /home/other\_user/}
\item PRM can block operations by detecting unexpected process ancestry combined
      with anomalous file access patterns
\item Existing MAC systems (e.g., SELinux, AppArmor) rely on static labels or
      executable paths, and lack mechanisms to reason about dynamic code provenance
      introduced via shared storage
\end{itemize}

\subsubsection{In-Process Exploitation via Malicious Python Payload}

\textbf{Attack}: In-Process Exploitation via Malicious Python Payload
\begin{itemize}
\item A Python-based HPC workload deserializes input data using pickle, permitting
      arbitrary code execution
\item Attacker injects a crafted pickle object via shared filesystem, causing
      malicious code execution within the same interpreter instance
\item Expected lineage: \texttt{slurm $\rightarrow$ user\_job $\rightarrow$ python $\rightarrow$ load\_model() $\rightarrow$ user\_data}
\item Observed lineage: \texttt{slurm $\rightarrow$ user\_job $\rightarrow$ python $\rightarrow$ pickle.load() $\rightarrow$ malicious\_payload $\rightarrow$ /home/other\_user/}
\item PRM can detect when injected code attempts actions inconsistent with the
      job's expected ancestry and semantic role
\item Existing MAC systems cannot distinguish this attack as execution remains
      within the same process, label, and credential context
\end{itemize}

\subsubsection{In-Process Code Execution via setup.py Abuse}

\textbf{Attack}: In-Process Code Execution via setup.py Abuse
\begin{itemize}
\item Python-based HPC workflows install dependencies at runtime using pip from
      shared filesystems
\item Attacker places malicious package in shared directory, causing arbitrary code
      execution during installation
\item Expected lineage: \texttt{slurm $\rightarrow$ user\_job $\rightarrow$ pip $\rightarrow$ python $\rightarrow$ setup.py $\rightarrow$ package\_install}
\item Observed lineage: \texttt{slurm $\rightarrow$ user\_job $\rightarrow$ pip $\rightarrow$ python $\rightarrow$ setup.py $\rightarrow$ malicious\_code}
\item PRM correlates runtime behavior with process ancestry to identify semantic
      deviations from expected installation logic
\item Existing MAC systems cannot detect this attack as malicious code executes
      within the same interpreter instance, under identical security labels and
      credentials
\end{itemize}

\subsubsection{Attempted Evasion via Ancestry Inflation}

\textbf{Attack}: Attempted Evasion via Ancestry Inflation
\begin{itemize}
\item Attacker attempts to evade PRM by artificially inflating process ancestry
      through repeated shell invocations
\item Job spawns a long chain of bash processes before executing malicious Python
      code to push original context beyond PRM's ancestry window
\item Expected lineage: \texttt{slurm $\rightarrow$ user\_job $\rightarrow$ python $\rightarrow$ legitimate\_computation}
\item Observed lineage: \texttt{slurm $\rightarrow$ user\_job $\rightarrow$ bash $\rightarrow$ bash $\rightarrow$ bash $\rightarrow$ python $\rightarrow$ malicious\_code}
\item PRM inspects last k ancestors (k=3), reducing observed ancestry to:
      \texttt{python $\rightarrow$ bash $\rightarrow$ bash}
\item Despite increased ancestry depth, semantic provenance remains unchanged: all
      processes execute under same job context and user identity
\item PRM detects malicious behavior when Python code performs actions inconsistent
      with expected semantic role, demonstrating that bounded ancestry inspection
      remains effective against evasion attempts
\end{itemize}

\subsubsection{Comprehensive Attack Coverage Analysis}

Beyond the specific scenarios detailed above, we evaluated PRM against a broader
range of HPC attack classes to assess comprehensive security coverage.
Table~\ref{tab:attack_coverage_comprehensive} summarizes PRM's effectiveness
across common HPC attack vectors.

\begin{table*}[htbp]
\centering
\caption{HPC Attack Classes Tested and Ongoing Evaluation by PRM}
\label{tab:attack_coverage_comprehensive}
\begin{tabular}{@{}p{4.5cm}p{4.5cm}p{2cm}@{}}
\toprule
\textbf{Attack Class} & \textbf{Example Techniques} & \textbf{PRM Coverage} \\
\midrule
Job impersonation & Stolen job scripts, wrapper abuse & Yes \\
Scheduler context abuse & Prolog/epilog manipulation & Yes \\
Shared filesystem data exfiltration & /scratch, /project abuse & Yes \\
Shared filesystem code injection & Script substitution, symlink attacks & Yes \\
Interpreter-level code execution & pickle, eval, setup.py & Yes \\
Dynamic runtime abuse & LD\_PRELOAD, env poisoning & Yes \\
Privilege abuse without UID change & Same-UID misuse, container root & Yes \\
setuid / sudo misuse & Policy-allowed escalation & Yes \\
Ancestry inflation (evasion attempt) & Long shell chains & Yes \\
Fork-based obfuscation & Process tree noise & Yes \\
\midrule
Kernel privilege escalation (post-exploit) & LPE followed by user-space abuse & Constrained \\
Root persistence mechanisms & Cron, systemd, init injection & Yes \\
Cross-user attack after root compromise & Accessing other users' data & Yes \\
Unauthorized kernel interface use & /proc, /sys, bpf, perf abuse & Yes \\
\midrule
Side-channel attacks & Cache timing, Spectre & No \\
Pure data poisoning & Training data corruption & No \\
Perfect semantic mimicry & Indistinguishable lineage and behavior & No \\
\bottomrule
\end{tabular}
\end{table*}

\subsubsection{Security Validation Results}

Through controlled laboratory testing, we evaluated PRM's approach to these
attack classes. The results indicate PRM's potential effectiveness in addressing
security challenges that are difficult to mitigate using traditional access
control mechanisms alone.

\subsection{Benchmark Methodology and Results}

\subsubsection{Test Environment}

Benchmarks were conducted on a deliberately resource-constrained Rocky Linux 9.5
VM to establish a conservative performance baseline:
\begin{itemize}
\item \textbf{Memory}: 2\,GB RAM
\item \textbf{Architecture}: x86\_64 (32-bit, 64-bit capable)
\item \textbf{Address Space}: 39 bits physical, 48 bits virtual
\item \textbf{CPUs}: 4 cores (0--3 online)
\item \textbf{Byte Order}: Little Endian
\end{itemize}

The minimal hardware configuration was chosen intentionally: because PRM's
per-syscall overhead is dominated by CPU-bound eBPF execution and hash
computation rather than memory bandwidth, measuring on a constrained platform
exposes worst-case latency that would be masked by the faster caches and higher
clock rates of production HPC nodes. Any overhead observed here therefore
represents an upper bound; on production hardware with higher single-thread
performance and larger CPU caches, absolute latencies are expected to decrease
proportionally. We acknowledge that validating this expectation on
production-scale HPC systems remains necessary future work
(Section~\ref{sec:evaluation}, Limitations).

\subsubsection{Benchmark Methodology}

We benchmark using 1,000,000 \texttt{open()-close()} syscalls to measure PRM's
overhead compared to baseline performance (the policy consists of two `allow'
system calls and one `deny' system call; denying a system call triggers an
exception that forces execution to skip subsequent system calls and jump directly
to the `catch' block). The test involves opening and immediately closing
\texttt{/tmp/test.dat} to isolate pure syscall overhead without I/O transfer
(access to the /tmp/ directory is restricted).

\textbf{Rationale for the \texttt{open/close} microbenchmark.}
File-open operations were selected because they are among the most
security-sensitive and performance-critical syscalls in HPC workloads: every
checkpoint write, shared-library load, log rotation, and data-stage operation
passes through the \texttt{file\_open} LSM hook. In PRM's architecture this hook
triggers the full five-layer processing pipeline --- cache lookup, fast-decision
filtering, ancestry extraction, rule evaluation, and PROG handler invocation ---
making it the heaviest per-call path and therefore the most informative single
syscall for comparing security framework overhead. Lighter hooks (e.g.,
\texttt{capable}) would understate overhead, while composite application
benchmarks would conflate security cost with I/O and compute variance. By
isolating the costliest hook we obtain a conservative, reproducible comparison
point between PRM and AppArmor under identical conditions. Broader workload
characterisation covering parallel I/O, MPI, and network-intensive patterns is
planned as future work on production HPC infrastructure.

\paragraph{Test Implementation}
\hfill
\begin{lstlisting}[language=Python]
def benchmark():
  filename = "/tmp/test.dat"
  iterations = 1000000
  
  for i in range(iterations):
    try:
      start_time = time.perf_counter()
      f = open(filename, "wb")
      f.close()
      end_time = time.perf_counter()
      
      if i % 1000 == 0:
        syscall_time = (end_time - start_time) * 1000000
        print(f"Iteration {i}: SUCCESS {syscall_time:.2f} us")
        
    except (PermissionError, OSError) as e:
      end_time = time.perf_counter()
      block_time = (end_time - start_time) * 1000000
      print(f"Iteration {i}: BLOCKED {block_time:.2f} us")
      
# Data collection used Linux pipeline: python3 benchmark.py | tee results.log
\end{lstlisting}

\begin{table}[htbp]
\centering
\caption{PRM vs AppArmor Performance Comparison (1M syscalls)}
\label{tab:benchmark_results}
\begin{tabular}{@{}p{2.8cm}|>{\columncolor{tableheader}}p{1.8cm}|p{1.8cm}|p{1.2cm}@{}}
\toprule
\rowcolor{tableheader}
\textbf{Configuration} & \textbf{Mean (\textmu s)} & \textbf{Std Dev (\textmu s)} & \textbf{Samples} \\
\midrule
Baseline (No Security) & 14.61 & 6.87  & 1,000 \\
\rowcolor{gray!10}
PRM Allow              & 23.36 & 14.90 & 1,000 \\
PRM Deny               & 20.39 & 9.66  & 1,000 \\
\rowcolor{gray!10}
AppArmor Allow         & 24.20 & 22.56 & 1,000 \\
AppArmor Block         & 12.84 & 10.78 & 1,000 \\
\bottomrule
\end{tabular}
\end{table}

\begin{figure}[htbp]
\centering
\includegraphics[width=0.9\columnwidth]{PRM_vs_AppArmor_LARGE.png}
\caption{Performance Distribution Analysis: Boxplot comparison of syscall latency
distributions for baseline, PRM allow, and PRM deny configurations. Outliers above
30\,\textmu s filtered for detailed view of core performance patterns.}
\label{fig:performance_measurements}
\end{figure}

\begin{table}[htbp]
\centering
\caption{Raw Performance Data Points (\textmu s per syscall)}
\label{tab:raw_data}
\begin{tabular}{@{}p{1.0cm}|>{\columncolor{tableheader}}p{1.0cm}|p{1.2cm}|p{1.0cm}|p{1.1cm}|p{1.1cm}@{}}
\toprule
\rowcolor{tableheader}
\textbf{Iteration} & \textbf{Baseline} & \textbf{PRM Allow} & \textbf{PRM Block} & \textbf{AppArmor Allow} & \textbf{AppArmor Block} \\
\midrule
0      & 49.59 & 186.01 & 104.63 & 73.44 & 38.12 \\
\rowcolor{gray!10}
100K   & 14.45 & 13.44  & 12.03  & 14.74 & 12.14 \\
200K   & 9.77  & 12.72  & 17.51  & 14.65 & 11.98 \\
\rowcolor{gray!10}
300K   & 22.34 & 27.04  & 23.95  & 14.60 & 12.09 \\
400K   & 12.99 & 19.11  & 35.71  & 14.76 & 11.99 \\
\rowcolor{gray!10}
500K   & 13.82 & 19.72  & 16.55  & 14.71 & 12.05 \\
600K   & 13.90 & 20.64  & 17.41  & 14.39 & 11.96 \\
\rowcolor{gray!10}
700K   & 14.05 & 18.02  & 16.80  & 20.19 & 12.03 \\
800K   & 13.31 & 33.47  & 16.68  & 81.47 & 12.03 \\
\rowcolor{gray!10}
900K   & 15.53 & 19.46  & 16.83  & 64.08 & 12.02 \\
\bottomrule
\end{tabular}
\end{table}

\subsubsection{Performance Analysis}

This section presents a direct syscall performance comparison between PRM's
eBPF-based LSM implementation and AppArmor's traditional LSM approach. AppArmor
provides standard path-based access control using pure LSM hooks for Debian
servers, while PRM offers advanced process ancestry analysis with intelligent
caching capabilities.

\begin{table}[htbp]
\centering
\caption{Syscall Performance Comparison: PRM vs AppArmor.
Formula: Average $= \sum(\text{sample\_times}) / \text{sample\_count}$}
\label{tab:syscall_comparison}
\begin{tabular}{@{}p{3cm}|>{\columncolor{tableheader}}p{2cm}|p{2.5cm}@{}}
\toprule
\rowcolor{tableheader}
\textbf{Configuration} & \textbf{Avg Latency (\textmu s)} & \textbf{Throughput (ops/sec)} \\
\midrule
Baseline (No Security) & 15.08 & 66,317 \\
\rowcolor{gray!10}
PRM Allow              & 23.62 & 42,334 \\
PRM Deny               & 20.67 & 48,385 \\
\rowcolor{gray!10}
AppArmor Allow         & 24.11 & 41,472 \\
AppArmor Block         & 12.84 & 77,880 \\
\bottomrule
\end{tabular}
\end{table}

\begin{table}[htbp]
\centering
\caption{Min/Max Latency Analysis Across Configurations}
\label{tab:minmax_analysis}
\begin{tabular}{@{}p{3cm}|>{\columncolor{tableheader}}p{2cm}|p{2cm}@{}}
\toprule
\rowcolor{tableheader}
\textbf{Configuration} & \textbf{Min (\textmu s)} & \textbf{Max (\textmu s)} \\
\midrule
Baseline (No Security) & 8.76  & 123.22 \\
\rowcolor{gray!10}
PRM Allow              & 11.79 & 271.58 \\
PRM Deny               & 10.94 & 111.56 \\
\rowcolor{gray!10}
AppArmor Allow         & 13.73 & 360.15 \\
AppArmor Block         & 11.74 & 348.15 \\
\bottomrule
\end{tabular}
\end{table}

The comparison demonstrates that PRM's ancestry analysis maintains competitive
per-call latency with AppArmor while providing security capabilities not
available in traditional LSM implementations.

\textbf{Contextualising the overhead.}
PRM Allow adds approximately 8.7\,\textmu s (\textasciitilde60\%) over the
no-security baseline. Three factors make this acceptable for HPC workloads:
\begin{enumerate}
\item \emph{Debug logging was enabled throughout all measurements.} Every syscall
      produced multiple \texttt{bpf\_trace\_printk} calls (see trace excerpt
      below), each of which acquires a global spinlock and writes to a ring
      buffer. Disabling debug output --- the expected production configuration ---
      removes this overhead entirely.
\item \emph{Caching amortises the cost.} The 8.7\,\textmu s figure reflects a
      cache-miss path. In a typical HPC loop performing $10^{6}$ identical
      operations, only the first invocation pays the full cost; the remaining
      $999{,}999$ resolve via an O(1) hash lookup whose latency is comparable to
      the baseline. The amortised per-call overhead therefore approaches zero as
      iteration count grows.
\item \emph{Syscall latency is a small fraction of total job runtime.} A
      scientific job running for hours accumulates syscall overhead measured in
      seconds, making the per-call difference negligible relative to compute time.
\end{enumerate}

\emph{Debug trace example (enabled during all benchmarks):}
\begin{lstlisting}[basicstyle=\scriptsize\ttfamily, frame=single, backgroundcolor=\color{gray!10}]
python3-3151 [003] ....1.1 5061.873207: bpf_trace_printk: Path: tmp
python3-3151 [003] ....1.1 5061.873302: bpf_trace_printk: INODE CREATE REJECTED
python3-3151 [003] ....1.1 5061.873329: bpf_trace_printk: PROG: Hook(9)-PID(3151)
{python3 -> bash -> sshd} Matched relation (REJECT)
\end{lstlisting}

The 3-level ancestry analysis (python3 --- bash --- sshd) demonstrates bounded
overhead due to the fixed-depth traversal and optimised rule table processing.

\subsection{Limitations and Future Work}

Current limitations include:
\begin{enumerate}
\item \textbf{Rule Complexity}: Manual rule creation requires security expertise
\item \textbf{Performance Validation}: Comprehensive benchmarks needed on
      production HPC systems
\item \textbf{Attack Coverage}: Security effectiveness needs validation against
      broader attack suites
\item \textbf{Management Tools}: GUI tools needed for rule creation and policy
      management
\item \textbf{Machine Learning}: Automated rule generation based on observed
      behavior patterns
\end{enumerate}

%%------------------------------------------------------------
%%  SECTION 5 — DISCUSSION
%%------------------------------------------------------------
\section{Discussion}
\label{sec:discussion}

\subsection{Novel Contributions}

The PRM Framework introduces several novel concepts not available in existing
security solutions:

\begin{enumerate}
\item \emph{Dynamic Rule Redirection}: Runtime policy adaptation through REDIRECT
      and RETURN actions --- a capability absent in static security solutions like
      SELinux and AppArmor
\item \emph{Intelligent Caching for Security}: Syscall-specific caching that
      eliminates redundant checks for repetitive HPC operations
\item \emph{Fast-Decision Filtering}: Pre-processing layer that bypasses expensive
      analysis for common patterns
\item \emph{Multi-Generational Ancestry Analysis}: Comprehensive process lineage
      tracking with configurable depth for context-aware security decisions
\item \emph{HPC-Optimized Security}: Specifically designed for the performance and
      security requirements of scientific computing environments
\end{enumerate}

\subsection{Impact on HPC Security Paradigm}

PRM represents a paradigm shift in HPC security:

\subsubsection{From Static to Dynamic Security}

Traditional HPC security relies on static policies that cannot adapt to execution
context. PRM introduces:
\begin{itemize}
\item Context-aware decisions: Same operation allowed/denied based on execution path
\item Runtime adaptation: Security policies that evolve with program execution
\item Intelligent caching: Performance optimization for repetitive scientific
      computations
\end{itemize}

\subsubsection{Enabling Safe User Code Execution}

HPC centers can now safely allow users to:
\begin{itemize}
\item Install custom software: Process ancestry ensures installations don't
      compromise system
\item Run privileged containers: Container operations monitored through ancestry
      analysis
\item Execute complex workflows: Multi-stage scientific pipelines secured through
      lineage tracking
\item Access shared resources: Fine-grained control over filesystem and network
      access
\end{itemize}

%%------------------------------------------------------------
%%  SECTION 6 — CONCLUSION
%%------------------------------------------------------------
\section{Conclusion}
\label{sec:conclusion}

The PRM Framework represents a significant advancement in HPC security by
introducing dynamic, context-aware security policies specifically designed for
scientific computing environments. Unlike existing static security solutions, PRM
adapts to execution context while maintaining the performance characteristics
essential for HPC workloads.

\subsection{Key Innovations}

\begin{enumerate}
\item Dynamic rule system: Runtime rule redirection and policy adaptation
      not available in existing static security solutions
\item Intelligent caching: Eliminates redundant security checks for repetitive
      operations (e.g., million-iteration syscalls like file writes, file opens,
      and network operations)
\item Fast-decision filtering: Pre-processing layer that bypasses expensive
      analysis for majority of syscalls
\item HPC user safety: Enables safe execution of user code while preventing
      system compromise
\item Multi-generational ancestry: Scalable process lineage analysis with
      configurable depth for context-aware security decisions
\end{enumerate}

\subsection{Future Directions}

Future work will focus on:
\begin{enumerate}
\item Comprehensive performance validation on production HPC systems, including
      parallel I/O, MPI, and network-intensive workloads, to confirm that the
      overhead upper bounds observed on our constrained test platform decrease on
      production hardware
\item Machine learning-based automated rule generation
\item Integration with HPC management and monitoring systems
\item Extension to distributed security policies across HPC clusters
\item Development of user-friendly policy management tools
\end{enumerate}

The PRM Framework opens new possibilities for secure, high-performance computing
by demonstrating that context-aware security and computational efficiency can
coexist in modern HPC environments.

%%------------------------------------------------------------
%%  DECLARATIONS
%%------------------------------------------------------------
\section*{Declarations}

\subsection*{Funding}
No funding was received for conducting this study.

\subsection*{Competing interests}
The authors have no competing interests to declare that are relevant to the
content of this article.

\subsection*{Data availability}
The data supporting the findings of this study are available from the corresponding
author upon reasonable request.

\subsection*{Author contributions}
Both authors contributed to the study conception, design, and implementation.
The first draft of the manuscript was written by Mojtaba Akbari and all authors
reviewed and approved the final manuscript.

%%------------------------------------------------------------
%%  REFERENCES
%%  Formatted per Springer Nature guidelines (numbered, square brackets)
%%------------------------------------------------------------
\bibliographystyle{sn-basic}
\begin{thebibliography}{99}

\bibitem{provos2003improving}
Provos N (2003) Improving host security with system call policies.
In: Proceedings of the 12th USENIX security symposium, USENIX Association

\bibitem{edge2015seccomp}
Edge J (2015) A seccomp overview. LWN.net

\bibitem{singh2019krsi}
Singh KP (2019) KRSI --- Kernel Runtime Security Instrumentation.
In: Linux security summit

\bibitem{salvan2016landlock}
Salvan M, Salaun M (2016) Landlock: unprivileged access control.
In: Linux security summit

\bibitem{google2018gvisor}
Google (2018) gVisor: container runtime sandbox.
\url{https://github.com/google/gvisor}

\bibitem{kata2018architecture}
Kata Containers (2018) Kata Containers architecture.
\url{https://katacontainers.io}

\bibitem{falco2016runtime}
Sysdig (2016) Falco: runtime security monitoring.
\url{https://falco.org}

\bibitem{cook2012yama}
Cook K (2012) YAMA Linux security module. Linux kernel documentation

\bibitem{li2018speaker}
Li Z et al (2018) SPEAKER: split-phase execution of application kernels for
runtime security. In: Proceedings of the 2018 ACM SIGSAC conference on computer
and communications security. ACM

\bibitem{hoiland2018express}
Hoiland-Jorgensen T et al (2018) The eXpress data path: fast programmable packet
processing in the operating system kernel. In: Proceedings of the 14th
international conference on emerging networking experiments and technologies
(CoNEXT). ACM

\bibitem{reshetova2016security}
Reshetova E, Karhunen J, Nyman T, Asokan N (2016) Security of OS-level
virtualization technologies. In: Nordic conference on secure IT systems. Springer,
Berlin, pp 77--93

\bibitem{xavier2013performance}
Xavier MG, Neves MV, Rossi FD, Ferreto TC, Lange T, De Rose CA (2013)
Performance evaluation of container-based virtualization for high performance
computing environments. In: 21st Euromicro international conference on parallel,
distributed, and network-based processing. IEEE, pp 233--240

\bibitem{nelson2017bpf}
Nelson L et al (2017) The BPF verifier. Linux kernel documentation

\bibitem{verizon2024dbir}
Verizon (2024) 2024 data breach investigations report (DBIR). Verizon

\bibitem{enisa2024threat}
ENISA (2024) Threat landscape 2024. European Union Agency for Cybersecurity

\bibitem{cisa2024vulns}
CISA (2024) Vulnerability bulletins and advisories 2024--2025. Weekly vulnerability
bulletins. Cybersecurity and Infrastructure Security Agency

\bibitem{microsoft2024rce}
Microsoft Security Blog (2024) CVE-2024-7971 zero-day RCE exploitation telemetry.
Microsoft

\bibitem{elastic2024pumakit}
Elastic Security Labs (2024) PUMAKIT Linux rootkit analysis. Elastic

\bibitem{fortinet2024lkm}
Fortinet (2024) LKM rootkit analysis and detection. Fortinet Blog

\bibitem{eset2024bootkitty}
ESET (2024) Bootkitty UEFI bootkit report. ESET Research

\bibitem{snyk2024container}
Snyk (2024) CVE-2024-21626 container escape analysis. Snyk Learn

\bibitem{vulncheck2024trends}
VulnCheck (2024) 2024 trends in vulnerability exploitation. VulnCheck

\bibitem{infosec2024kevs}
Infosecurity Magazine (2024) Zero-day and known exploited vulnerabilities report.
Infosecurity Magazine

\bibitem{globenews2024dbir}
GlobeNewswire (2024) 2024 Verizon data breach investigations report press release.
GlobeNewswire

\bibitem{hackernews2024cves}
The Hacker News (2024) CVE exploitation statistics coverage. The Hacker News

\bibitem{crowdstrike2024gtr}
CrowdStrike (2024) Global threat report 2024. CrowdStrike

\bibitem{mandiant2024mtrends}
Mandiant (2024) M-Trends 2024. Mandiant

\bibitem{mitre2024attack}
MITRE (2024) ATT\&CK prevalence report 2024. MITRE Corporation

\bibitem{wiz2024cloud}
Wiz (2024) Cloud threat report 2024. Wiz

\bibitem{owasp2024top10}
OWASP (2024) Top 10 web application security risks 2024. OWASP Foundation

\bibitem{zhang2021analyzing}
Zhang L, Jaeger T, Liu P (2021) Analyzing the overhead of filesystem protection
using Linux security modules. In: Proceedings of the 2021 ACM SIGSAC conference
on computer and communications security. ACM

\bibitem{pappas2013kbouncer}
Pappas V, Polychronakis M, Keromytis AD (2013) Transparent ROP exploit mitigation
using indirect branch tracing. In: Proceedings of the 22nd USENIX security
symposium. USENIX Association, pp 447--462

\bibitem{ghavamnia2020confine}
Ghavamnia H et al (2020) Confine: automated system call policy generation for
container attack surface reduction. In: Proceedings of the 23rd international
symposium on research in attacks, intrusions and defenses (RAID), pp 443--458

\bibitem{wang2021semantic}
Wang R et al (2021) Towards semantic-aware security monitoring for containerized
applications. In: Proceedings of the 42nd IEEE symposium on security and privacy
(Oakland). IEEE, pp 1352--1369

\bibitem{vasilakis2019hpc}
Vasilakis N et al (2019) Security of high performance computing systems: a
comprehensive survey. ACM Comput Surv 52(6):1--36

\bibitem{nelson2020ebpf}
Nelson L et al (2020) eBPF-based security enforcement with minimal overhead.
In: Proceedings of the 2020 USENIX annual technical conference (ATC).
USENIX Association, pp 615--628

\bibitem{gerhardt2017shifter}
Gerhardt L et al (2017) Shifter: containers for HPC. J Phys Conf Ser 898.
IOP Publishing

\bibitem{priedhorsky2017charliecloud}
Priedhorsky R, Randles T (2017) Charliecloud: unprivileged containers for
user-defined software stacks in HPC. In: Proceedings of the international
conference for high performance computing, networking, storage and analysis (SC)

\bibitem{younge2017tale}
Younge AJ et al (2017) A tale of two systems: using containers to deploy HPC
applications on supercomputers and clouds. In: 2017 IEEE international conference
on cloud computing technology and science (CloudCom). IEEE, pp 74--81

\bibitem{jacobsen2015contain}
Jacobsen DM, Canon RS (2015) Contain this, unleashing Docker for HPC.
In: Proceedings of the Cray user group (CUG)

\bibitem{canon2016scream}
Canon RS, Jacobsen DM (2016) Scream: container-based resource management for HPC
workloads. In: 2016 IEEE international conference on cluster computing (CLUSTER).
IEEE, pp 7--15

\bibitem{kurtzer2017singularity}
Kurtzer GM, Sochat V, Bauer MW (2017) Singularity: scientific containers for
mobility of compute. PLoS ONE 12(5):e0177459

\bibitem{rudyy2018containers}
Rudyy O et al (2018) Containers in HPC: a scalability and portability study in
production biological simulations. In: 2018 IEEE international parallel and
distributed processing symposium (IPDPS). IEEE, pp 567--577

\bibitem{beltre2019enabling}
Beltre AM et al (2019) Enabling HPC workloads on cloud infrastructure using
Kubernetes container orchestration mechanisms. In: 2019 IEEE international
parallel and distributed processing symposium workshops (IPDPSW). IEEE, pp 11--20

\bibitem{torrez2019hpc}
Torrez A et al (2019) HPC container runtimes have minimal or no performance
impact. In: 2019 IEEE/ACM international workshop on containers and new
orchestration paradigms for isolated environments in HPC (CANOPIE-HPC).
IEEE, pp 37--42

\bibitem{higgins2017orchestrating}
Higgins J et al (2017) Orchestrating Docker containers in the HPC environment
with Shifter. In: 2017 international conference on high performance computing
\& simulation (HPCS). IEEE, pp 506--513

\end{thebibliography}

\end{document}
